% =====================================================================
\chapter{Probabilidad condicionada, independencia y teorema de Bayes}
\label{tema:4}
% =====================================================================

La probabilidad condicionada formaliza cómo cambia la incertidumbre al recibir información. De ella se deducen la regla de la cadena, el teorema de la probabilidad total y el teorema de Bayes, y con ella se define la independencia. El tema sigue las secciones §1.3 a §1.5 de Bertsekas y Tsitsiklis~\cite{bt} (clases~2 y~3 de MIT 6.041~\cite{mit6041f10}), el capítulo~2 de Blitzstein y Hwang~\cite{bh}, el capítulo~3 de Ross~\cite{ross} y el capítulo~17 de~\cite{mcs}, cuyo tratamiento del problema de Monty Hall y de las pruebas médicas (§17.1 y §17.4.2) se adapta aquí.

\section{Probabilidad condicionada}

\begin{definicion}[Probabilidad condicionada]
Sean $A$ y $B$ sucesos con $\PP(B)>0$. La \emph{probabilidad de $A$ condicionada a $B$} es
\[
\PP(A\mid B)=\frac{\PP(A\cap B)}{\PP(B)}.
\]
\end{definicion}

La idea es que, al saber que ha ocurrido $B$, los resultados fuera de $B$ quedan descartados y los de $B$ conservan sus probabilidades \emph{relativas}; dividir por $\PP(B)$ restablece la normalización. En el modelo de Laplace, $\PP(A\mid B)=\abs{A\cap B}/\abs{B}$: se cuenta dentro de $B$ como si $B$ fuera el nuevo espacio muestral.

\begin{proposicion}\label{prop:t4-esprob}
Si $\PP(B)>0$, la función $A\mapsto\PP(A\mid B)$ es una ley de probabilidad sobre $\Omega$. En consecuencia, cumple todas las propiedades de la \cref{prop:t3-props}; por ejemplo, $\PP(A^{c}\mid B)=1-\PP(A\mid B)$.
\end{proposicion}

\begin{proof}
No negatividad: $\PP(A\cap B)\ge0$ y $\PP(B)>0$. Normalización: $\PP(\Omega\mid B)=\PP(B)/\PP(B)=1$. Aditividad: si $A_1,A_2,\dots$ son incompatibles, también lo son $A_i\cap B$, y $\bigl(\bigcup_iA_i\bigr)\cap B=\bigcup_i(A_i\cap B)$; por el axioma~(iii) aplicado a $\PP$, $\PP\bigl(\bigcup_iA_i\mid B\bigr)=\sum_i\PP(A_i\cap B)/\PP(B)=\sum_i\PP(A_i\mid B)$.
\end{proof}

\begin{ejemplo}\label{ej:t4-hijos}
Una familia tiene dos hijos; los cuatro casos (chico o chica para el mayor y para el menor) son equiprobables. Si se sabe que el mayor es chica, la probabilidad de que ambos lo sean es $\frac{1/4}{1/2}=\frac12$. Si solo se sabe que \emph{al menos uno} es chica, el suceso condicionante tiene tres casos y la probabilidad es $\frac{1/4}{3/4}=\frac13$. Informaciones de apariencia parecida condicionan a sucesos distintos.
\end{ejemplo}

\section{Regla de la cadena}

\begin{teorema}[Regla de la cadena o de la multiplicación]\label{teo:t4-cadena}
Si $A_1,\dots,A_n$ son sucesos con $\PP(A_1\cap\dots\cap A_{n-1})>0$, entonces
\[
\PP(A_1\cap A_2\cap\dots\cap A_n)=\PP(A_1)\,\PP(A_2\mid A_1)\,\PP(A_3\mid A_1\cap A_2)\cdots\PP(A_n\mid A_1\cap\dots\cap A_{n-1}).
\]
\end{teorema}

\begin{proof}
Todas las probabilidades condicionadas están definidas, porque $A_1\cap\dots\cap A_k\supseteq A_1\cap\dots\cap A_{n-1}$ tiene probabilidad positiva para $k\le n-1$. Por inducción sobre $n$: para $n=1$ no hay nada que probar, y si la fórmula vale para $n-1$, la definición de probabilidad condicionada con $B=A_1\cap\dots\cap A_{n-1}$ da $\PP(B\cap A_n)=\PP(B)\,\PP(A_n\mid B)$; basta sustituir $\PP(B)$ por la fórmula para $n-1$ sucesos.
\end{proof}

\begin{ejemplo}
Se extraen tres cartas sin reemplazamiento. Si $A_i$ es «la carta $i$ es de corazones», $\PP(A_1\cap A_2\cap A_3)=\frac{13}{52}\cdot\frac{12}{51}\cdot\frac{11}{50}\approx0{,}0129$. La regla de la cadena sustituye un recuento global por una sucesión de razonamientos locales, que es lo que se representa en un \emph{diagrama de árbol}: la probabilidad de cada hoja es el producto de las probabilidades de las ramas que llevan a ella~\cite[§16.2]{mcs}.
\end{ejemplo}

\section{Probabilidad total y teorema de Bayes}

\begin{definicion}
Una familia de sucesos $B_1,\dots,B_n$ es una \emph{partición} de $\Omega$ si son incompatibles dos a dos y $B_1\cup\dots\cup B_n=\Omega$.
\end{definicion}

\begin{teorema}[Probabilidad total]\label{teo:t4-total}
Si $B_1,\dots,B_n$ es una partición de $\Omega$ con $\PP(B_i)>0$ para todo $i$, entonces, para todo suceso $A$,
\[
\PP(A)=\sum_{i=1}^{n}\PP(A\mid B_i)\,\PP(B_i).
\]
\end{teorema}

\begin{proof}
$A=\bigcup_i(A\cap B_i)$ y los sucesos $A\cap B_i$ son incompatibles dos a dos. Por la aditividad, $\PP(A)=\sum_i\PP(A\cap B_i)$, y $\PP(A\cap B_i)=\PP(A\mid B_i)\PP(B_i)$ por definición.
\end{proof}

\begin{teorema}[Bayes]\label{teo:t4-bayes}
En las condiciones del \cref{teo:t4-total}, si además $\PP(A)>0$,
\[
\PP(B_j\mid A)=\frac{\PP(A\mid B_j)\,\PP(B_j)}{\sum_{i=1}^{n}\PP(A\mid B_i)\,\PP(B_i)}.
\]
\end{teorema}

\begin{proof}
$\PP(B_j\mid A)=\PP(A\cap B_j)/\PP(A)$; el numerador es $\PP(A\mid B_j)\PP(B_j)$ y el denominador se desarrolla con el \cref{teo:t4-total}.
\end{proof}

En la terminología bayesiana, $\PP(B_j)$ es la probabilidad \emph{a priori} de la hipótesis $B_j$, $\PP(A\mid B_j)$ es la \emph{verosimilitud} de la evidencia $A$ bajo esa hipótesis y $\PP(B_j\mid A)$ es la probabilidad \emph{a posteriori}. Con dos hipótesis $B$ y $B^{c}$, dividiendo las dos fórmulas de Bayes se obtiene la \emph{forma de razones} (\emph{odds}):
\[
\frac{\PP(B\mid A)}{\PP(B^{c}\mid A)}=\frac{\PP(B)}{\PP(B^{c})}\cdot\frac{\PP(A\mid B)}{\PP(A\mid B^{c})},
\]
«razón a posteriori = razón a priori $\times$ razón de verosimilitudes».

\begin{ejemplo}[Prueba diagnóstica]\label{ej:t4-diagnostico}
Una enfermedad afecta al $1\,\%$ de la población. Una prueba detecta el $90\,\%$ de los enfermos (\emph{sensibilidad}) y da positivo en el $5\,\%$ de los sanos (tasa de falsos positivos; la \emph{especificidad} es el $95\,\%$). Sean $E$ «estar enfermo» y $+$ «dar positivo». Por Bayes,
\[
\PP(E\mid +)=\frac{0{,}9\cdot0{,}01}{0{,}9\cdot0{,}01+0{,}05\cdot0{,}99}=\frac{0{,}009}{0{,}0585}=\frac{2}{13}\approx0{,}154.
\]
Pese a la aparente calidad de la prueba, solo un $15\,\%$ de los positivos están enfermos. En \emph{frecuencias naturales}: de cada 1000 personas, 10 están enfermas y 9 dan positivo; de las 990 sanas, unas 49,5 dan positivo; de unos 58,5 positivos, solo 9 son enfermos~\cite[§17.4.2 y §17.4.4]{mcs}.
\end{ejemplo}

\begin{ejemplo}[Monty Hall]\label{ej:t4-monty}
Un concursante elige una de tres puertas; tras una hay un coche y tras las otras dos, cabras. El presentador, que sabe dónde está el coche, abre siempre una de las otras dos puertas que oculte una cabra (al azar si hay dos) y ofrece cambiar. Sin pérdida de generalidad, el concursante elige la puerta~1. Sea $C_i$ «el coche está en la puerta $i$», con $\PP(C_i)=1/3$, y $M_3$ «el presentador abre la puerta~3». Entonces $\PP(M_3\mid C_1)=1/2$, $\PP(M_3\mid C_2)=1$ y $\PP(M_3\mid C_3)=0$. Por Bayes,
\[
\PP(C_2\mid M_3)=\frac{1\cdot\frac13}{\frac12\cdot\frac13+1\cdot\frac13+0\cdot\frac13}=\frac23,\qquad\PP(C_1\mid M_3)=\frac13.
\]
Cambiar de puerta duplica la probabilidad de ganar. La intuición «quedan dos puertas, luego $1/2$» ignora que la elección del presentador \emph{depende} de dónde está el coche y, por tanto, aporta información~\cite[§16.1 y §17.1]{mcs}.
\end{ejemplo}

\section{Independencia}

\begin{definicion}[Independencia]
Dos sucesos $A$ y $B$ son \emph{independientes} si $\PP(A\cap B)=\PP(A)\PP(B)$. Si $\PP(B)>0$, esto equivale a $\PP(A\mid B)=\PP(A)$: saber que ha ocurrido $B$ no cambia la probabilidad de $A$. Los sucesos $A_1,\dots,A_n$ son \emph{independientes} (o mutuamente independientes) si para todo subconjunto de índices $S\subseteq[n]$ con $\abs{S}\ge2$,
\[
\PP\Bigl(\bigcap_{i\in S}A_i\Bigr)=\prod_{i\in S}\PP(A_i).
\]
Son \emph{independientes dos a dos} si la condición se cumple al menos para todos los $S$ con $\abs{S}=2$.
\end{definicion}

\begin{proposicion}\label{prop:t4-complementos}
Si $A$ y $B$ son independientes, también lo son $A$ y $B^{c}$, $A^{c}$ y $B$, y $A^{c}$ y $B^{c}$.
\end{proposicion}

\begin{proof}
$A$ es la unión disjunta de $A\cap B$ y $A\cap B^{c}$, luego $\PP(A\cap B^{c})=\PP(A)-\PP(A)\PP(B)=\PP(A)(1-\PP(B))=\PP(A)\PP(B^{c})$. El segundo caso es simétrico, y el tercero resulta de aplicar el primero a la pareja independiente $A^{c}$, $B$.
\end{proof}

\begin{ejemplo}[Independencia dos a dos sin independencia mutua]\label{ej:t4-dosados}
Se lanzan dos monedas equilibradas; sean $A$ «la primera es cara», $B$ «la segunda es cara» y $C$ «las dos coinciden». Cada suceso tiene probabilidad $1/2$ y cada intersección de dos tiene probabilidad $1/4$ (por ejemplo, $A\cap C$ es «dos caras»), luego son independientes dos a dos. Pero $\PP(A\cap B\cap C)=1/4\neq1/8$: conocer $A$ y $B$ determina $C$.
\end{ejemplo}

\begin{observacion}[Independencia e incompatibilidad]
No deben confundirse. Si $A$ y $B$ son incompatibles con probabilidades positivas, \emph{no} son independientes: $\PP(A\cap B)=0<\PP(A)\PP(B)$; saber que ocurre $A$ garantiza que no ocurre $B$. La independencia es, además, una propiedad de la ley $\PP$, no de los sucesos como conjuntos: en la práctica suele ser una \emph{hipótesis} de modelización (lanzamientos sucesivos, ensayos repetidos) más que algo que se demuestre~\cite[§17.7.2]{mcs}.
\end{observacion}

Si $\PP(C)>0$, los sucesos $A$ y $B$ son \emph{condicionalmente independientes dado $C$} si $\PP(A\cap B\mid C)=\PP(A\mid C)\PP(B\mid C)$. Ni la independencia implica la independencia condicional ni al revés (\cref{eje:4-ajedrez}).

\begin{observacion}[Conexión con el curso]
La regla de la cadena es el fundamento de los modelos de lenguaje autorregresivos (módulo~1): $\PP(x_1,\dots,x_T)=\prod_{t=1}^{T}\PP(x_t\mid x_1,\dots,x_{t-1})$, y un Transformer (módulos~13 a~15) calcula cada factor. La hipótesis de Markov del modelo de $n$-gramas es una hipótesis de independencia condicional. El teorema de Bayes reaparece en el módulo~5 (inferencia bayesiana) y en los clasificadores del módulo~8.
\end{observacion}

\begin{observacion}[Errores frecuentes]
(1) Confundir $\PP(A\mid B)$ con $\PP(B\mid A)$ (la \emph{falacia del fiscal}): que una prueba detecte al $90\,\%$ de los enfermos no significa que el $90\,\%$ de los positivos esté enfermo (\cref{ej:t4-diagnostico}). (2) Olvidar la probabilidad a priori, que en enfermedades raras domina el resultado. (3) Confundir independencia con incompatibilidad. (4) Deducir la independencia mutua de la independencia dos a dos (\cref{ej:t4-dosados}). (5) Condicionar al suceso equivocado: «al menos uno es chica» no es lo mismo que «el mayor es chica» (\cref{ej:t4-hijos}).
\end{observacion}

\begin{lecturas}
\item Bertsekas y Tsitsiklis~\cite{bt}, §1.3 a §1.5 (clases~2 y~3 de MIT 6.041~\cite{mit6041f10}).
\item Blitzstein y Hwang~\cite{bh}, cap.~2 «Conditional Probability», en especial §2.3 «Bayes' rule and the law of total probability», §2.5 «Independence of events» y §2.8 «Pitfalls and paradoxes».
\item Ross~\cite{ross}, cap.~3, §3.2 «Conditional Probabilities», §3.3 «Bayes's Formula» y §3.4 «Independent Events».
\item Lehman, Leighton y Meyer~\cite{mcs}, cap.~17 «Conditional Probability», §17.1 a §17.5, §17.7 y §17.8.
\item Grinstead y Snell~\cite{gs}, §4.1 «Discrete Conditional Probability» y §4.3 «Paradoxes».
\end{lecturas}

\section*{Ejercicios del tema 4}
\addcontentsline{toc}{section}{Ejercicios del tema 4}

\begin{ejercicio}{1}{eje:4-tabla}
En una clase de 100 estudiantes, 40 son de Ingeniería y 60 de Matemáticas; de los de Ingeniería, 30 aprueban, y de los de Matemáticas, 36. Se elige un estudiante al azar. Calcule $\PP(\text{aprueba})$, $\PP(\text{aprueba}\mid\text{Ingeniería})$ y $\PP(\text{Ingeniería}\mid\text{aprueba})$. ¿Son independientes «aprobar» y «ser de Ingeniería»?
\end{ejercicio}

\begin{ejercicio}{2}{eje:4-cachorros}
(a) Una perra tiene dos cachorros; se sabe que el mayor es hembra. ¿Qué probabilidad hay de que los dos lo sean? (b) Una gata tiene dos gatitos; se sabe que al menos uno es macho. ¿Qué probabilidad hay de que los dos lo sean? Suponga sexos equiprobables e independientes.
\fuente{adaptado de MIT 18.05~\cite{mit1805}, \emph{Problem Set}~2, problema~1.}
\end{ejercicio}

\begin{ejercicio}{2}{eje:4-spam}
El $80\,\%$ de los correos que recibe una cuenta es correo basura. La expresión «dinero gratis» aparece en el $10\,\%$ de los correos basura y en el $1\,\%$ de los legítimos. Si un correo contiene esa expresión, ¿qué probabilidad hay de que sea basura?
\fuente{adaptado de Harvard Stat 110~\cite{stat110}, \emph{Strategic Practice}~2, sección «Thinking Conditionally», problema~2.}
\end{ejercicio}

\begin{ejercicio}{2}{eje:4-dos-pruebas}
Una enfermedad afecta al $1\,\%$ de la población y una prueba tiene sensibilidad y especificidad del $95\,\%$. Una persona da positivo en dos pruebas que, condicionadas al estado de salud, son independientes. Calcule la probabilidad de que esté enferma (a) actualizando con las dos pruebas a la vez; (b) actualizando primero con una y después con la otra. Compruebe que coinciden.
\fuente{adaptado de Harvard Stat 110~\cite{stat110}, \emph{Homework}~3, problema~2(b).}
\end{ejercicio}

\begin{ejercicio}{2}{eje:4-piezas}
Un almacén tiene 500 piezas antiguas, de las que el $15\,\%$ son defectuosas, y 1500 nuevas, con un $5\,\%$ de defectuosas. Se lanza una moneda para elegir el tipo y se toman dos piezas de ese tipo, con reemplazamiento. (a) ¿Qué probabilidad hay de que ambas sean defectuosas? (b) Si ambas lo son, ¿qué probabilidad hay de que fueran antiguas?
\fuente{adaptado de MIT 6.041~\cite{mit6041f10}, \emph{Problem Set}~2, problema~3 (se ha fijado el muestreo con reemplazamiento).}
\end{ejercicio}

\begin{ejercicio}{2}{eje:4-perro}
Un perro se ha perdido en el bosque A con probabilidad $0{,}4$ o en el B con probabilidad $0{,}6$. Si está en A y se busca allí un día, se encuentra con probabilidad $0{,}25$; si está en B y se busca allí, con probabilidad $0{,}15$. (a) ¿Dónde conviene buscar el primer día? (b) Si se busca en A y no se encuentra, ¿qué probabilidad hay de que esté en A? (c) Si se elige el bosque lanzando una moneda y el perro aparece el primer día, ¿qué probabilidad hay de que se buscara en A?
\fuente{adaptado de MIT 6.041~\cite{mit6041f10}, \emph{Problem Set}~2, problema~4, apartados (a) a~(c).}
\end{ejercicio}

\begin{ejercicio}{2}{eje:4-monty7}
(a) En una versión del concurso con 7 puertas, el concursante elige una y el presentador, que sabe dónde está el coche, abre 3 de las otras 6, todas con cabra. Si el concursante cambia a una de las 3 puertas cerradas restantes elegida al azar, ¿qué probabilidad tiene de ganar? ¿Le conviene cambiar? (b) Generalice a $n\ge3$ puertas, de las que el presentador abre $m$, con $1\le m\le n-2$.
\fuente{adaptado de Harvard Stat 110~\cite{stat110}, \emph{Homework}~3, problema~1.}
\end{ejercicio}

\begin{ejercicio}{2}{eje:4-monty-p}
En el problema de Monty Hall (\cref{ej:t4-monty}), supóngase que, cuando el presentador puede abrir la puerta~2 o la~3, abre la~2 con probabilidad $p\in[\tfrac12,1]$. Calcule la probabilidad de ganar cambiando (a) sin condicionar; (b) sabiendo que abre la puerta~2; (c) sabiendo que abre la puerta~3.
\fuente{adaptado de Harvard Stat 110~\cite{stat110}, \emph{Strategic Practice}~3, sección «Continuing with Conditioning», problema~1.}
\end{ejercicio}

\begin{ejercicio}{2}{eje:4-paridad}
Se lanzan dos dados. Sean $A$ «el primero es impar», $B$ «el segundo es impar» y $C$ «la suma es impar». Estudie si son independientes dos a dos y si son mutuamente independientes.
\fuente{adaptado de MIT 18.05~\cite{mit1805}, \emph{Problem Set}~3, problema~1(a).}
\end{ejercicio}

\begin{ejercicio}{2}{eje:4-ajedrez}
Un jugador se enfrenta a un rival desconocido que, con probabilidad $1/3$ cada caso, es principiante, intermedio o maestro; contra ellos gana cada partida con probabilidad $0{,}9$, $0{,}5$ y $0{,}3$, respectivamente, y las partidas son independientes \emph{dado} el nivel del rival. (a) Calcule la probabilidad de ganar la primera partida. (b) Calcule la probabilidad de ganar la segunda sabiendo que ganó la primera. (c) ¿Son independientes los resultados de las dos partidas? Explique la diferencia entre independencia e independencia condicional.
\fuente{adaptado de Harvard Stat 110~\cite{stat110}, \emph{Strategic Practice}~2, sección «Thinking Conditionally», problema~5.}
\end{ejercicio}

\begin{ejercicio}{3}{eje:4-si-mismo}
Demuestre que un suceso $A$ es independiente de sí mismo si y solo si $\PP(A)\in\{0,1\}$. Demuestre además que un suceso con $\PP(A)\in\{0,1\}$ es independiente de cualquier otro suceso.
\fuente{adaptado de Harvard Stat 110~\cite{stat110}, \emph{Strategic Practice}~2, sección «Independence», problema~1.}
\end{ejercicio}

\begin{ejercicio}{3}{eje:4-union-indep}
(a) Demuestre que si $A$, $B$ y $C$ son mutuamente independientes, entonces $A$ y $B\cup C$ son independientes. (b) Demuestre que si $A_1,\dots,A_n$ son independientes, $\PP\bigl(\bigcup_{i=1}^{n}A_i\bigr)=1-\prod_{i=1}^{n}\bigl(1-\PP(A_i)\bigr)$.
\fuente{adaptado de MIT 6.041~\cite{mit6041f10}, \emph{Problem Set}~2, problema~5.}
\end{ejercicio}

\begin{ejercicio}{3}{eje:4-producto-triple}
Construya tres sucesos $A$, $B$ y $C$ que cumplan $\PP(A\cap B\cap C)=\PP(A)\PP(B)\PP(C)$ pero que no sean mutuamente independientes, y demuestre que su ejemplo cumple ambas propiedades.
\fuente{adaptado de Harvard Stat 110~\cite{stat110}, \emph{Strategic Practice}~2, sección «Independence», problema~4.}
\end{ejercicio}

\begin{ejercicio}{3}{eje:4-razones}
(a) Demuestre la forma de razones del teorema de Bayes. (b) Sean $E_1$ y $E_2$ dos evidencias condicionalmente independientes dadas $B$ y dadas $B^{c}$. Demuestre que actualizar la razón a priori con $E_1$ y después con $E_2$ da el mismo resultado que actualizar con $E_1\cap E_2$ de una vez.
\fuente{el apartado (a) está adaptado de Harvard Stat 110~\cite{stat110}, \emph{Homework}~3, problema~2(a).}
\end{ejercicio}

\begin{ejercicio}{3}{eje:4-simetria}
Sean $A$ y $B$ sucesos con probabilidades estrictamente entre 0 y 1. Demuestre que las tres afirmaciones siguientes son equivalentes: (a) $\PP(A\mid B)>\PP(A)$; (b) $\PP(B\mid A)>\PP(B)$; (c) $\PP(A\mid B)>\PP(A\mid B^{c})$.
\end{ejercicio}
